\documentclass[../main.tex]{subfiles}
\graphicspath{{\subfix{../image/}}{\subfix{../figure/}}}

\begin{document}

\chapter*{引言}\label{chapter:introduction}
\addcontentsline{toc}{chapter}{引言}

测度值解在变分学与非线性偏微分方程（PDE）问题中扮演重要角色，而这类问题源自物理、化学、生物与工程。

谈论“测度值解”其实是有歧义的，因为这一术语同时指代 Radon 测度值解与 Young 测度值解这两种概念；它们却惊人地契合于同一数学框架。下面扼要列举若干研究领域，其中这两种概念之一或两者都发挥过重要作用，这便是本节的内容。不过本节并不打算涵盖测度值解所出现的全部发展问题谱（特别地，关于概率测度的 Wasserstein 空间中梯度流理论及其在输运方程上的应用，我们请读者参见文献~\cite{ref06}）。

\section*{动机}\label{section:intro-motivation}

\subsection*{椭圆与抛物 PDE 的 Radon 测度值解}\label{subsection:intro-radon}

在椭圆情形，我们关心形如
\begin{align*}
Lu = \mu \quad \text{在 }\Omega, \qquad(E0)
\end{align*}
的方程，其中 $\Omega\subseteq\mathbb{R}^N$ 为开集，$u=u(x)$（$x\in\Omega$），$L$ 为线性或非线性椭圆算子，$\mu$ 为 $\Omega$ 上的 Radon 测度。若 $\Omega\subsetneq\mathbb{R}^N$，则在边界 $\partial\Omega$ 上须满足边界条件
\begin{align*}
u=\nu \quad \text{在 }\partial\Omega, \qquad(BC)
\end{align*}
边界数据 $\nu$ 本身也可以是 Radon 测度。

在抛物情形，我们考虑形如
\begin{align*}
\begin{cases}
\partial_t u + Lu = \mu & \text{在 }\Omega\times(0,T)=:Q,\\
u=u_0 & \text{在 }\Omega\times\{0\},
\end{cases}
\qquad(P0)
\end{align*}
的初值问题，其中 $\Omega$ 同上，$T>0$，$u=u(x,t)$（$(x,t)\in Q$），且 $\mu$ 与 $u_0$ 中至少一个为 Radon 测度。同样，若 $\Omega\subsetneq\mathbb{R}^N$，则在侧边界 $\partial\Omega\times(0,T)$ 上须施加合适的边界条件，其中同样可涉及 Radon 测度。

方程(E0)与问题(P0)的一个范例分别是 Poisson 方程与热方程的 Cauchy 问题。这两个情形（对应于取 $\Omega=\mathbb{R}^N$、$L=-\Delta$、在(E0)中 $\mu=\delta_{x_0}$，以及在(P0)中 $\mu=0$、$u_0=\delta_{x_0}$，其中 $\delta_{x_0}$ 为集中于某点 $x_0\in\Omega$ 的 Dirac 质量）在线性 PDE 理论中占据中心地位，并引导人们研究基本解；由线性性对基本解作叠加，即可表示更一般的解。

在线性理论的基础上，人们自然希望在非线性情形研究类似问题。然而很快就发现，即使在简单的半线性情形，在通常的 $L^1$ 框架内解的存在性本身也会遇到障碍。

\subsubsection*{椭圆情形}\label{subsubsection:intro-elliptic}

文献~\cite{ref14}以生动的方式描述了第一个不存在性案例，它涉及半线性椭圆方程
\begin{align}
-\Delta u + |u|^{p-1}u = \mu \quad \text{在 }\mathbb{R}^N\ (p>0), \label{align:intro-0.1}
\end{align}
其中 $N\geqslant 3$。方程\eqref{align:intro-0.1}与方程
\begin{align}
\lambda v - \Delta(|v|^{\alpha-1}v) = \mu \quad \text{在 }\mathbb{R}^N\ (\lambda,\alpha>0) \label{align:intro-0.2}
\end{align}
相关联（取 $u=|v|^{\alpha-1}v$、$p=1/\alpha$，并缩放掉 $\lambda$），而后者正是多孔介质方程
\begin{align}
\partial_t v = \Delta(|v|^{\alpha-1}v) \quad \text{在 }\mathbb{R}^N\times(0,T)\ (\alpha,T>0) \label{align:intro-0.3}
\end{align}
的 resolvent 方程。既然非线性半群理论建议在函数空间 $L^1(\mathbb{R}^N)$ 中研究方程\eqref{align:intro-0.3}，那么在\eqref{align:intro-0.1}–\eqref{align:intro-0.2}中自然应假设 $\mu\in L^1(\mathbb{R}^N)$。研究\eqref{align:intro-0.1}的另一个动机来自大原子电子密度的 Thomas--Fermi 模型，此时 $\mu$ 是 Dirac 质量的有限叠加~\cite{ref65}。

设 $\Omega\subseteq\mathbb{R}^N$ 为具有光滑边界 $\partial\Omega$ 的有界开集。文献~\cite{ref28}证明，对任意 $\mu\in L^1(\Omega)$ 与 $p>0$，齐次 Dirichlet 问题
\[
\begin{cases}
-\Delta u + |u|^{p-1}u = \mu & \text{在 }\Omega,\\
u=0 & \text{在 }\partial\Omega
\end{cases}
\]
存在唯一解 $u\in W^{1,1}_0(\Omega)$，且 $|u|^p\in L^1(\Omega)$。方程\eqref{align:intro-0.1}的适定性结果随后在文献~\cite{ref15}中给出。

关于\eqref{align:intro-0.1}右端为 Radon 测度的情形，人们很快认识到条件
\begin{align}
p < \frac{N}{N-2} \label{align:intro-0.4}
\end{align}
起着核心作用。事实上，记 $\mathscr{R}_f(\Omega)$ 为 $\Omega\subseteq\mathbb{R}^N$ 上有限（带号）Radon 测度所成的 Banach 空间，赋以范数 $\|\mu\|_{\mathscr{R}_f(\Omega)}:=|\mu|(\Omega)$（$|\mu|(\Omega)<\infty$ 为测度 $\mu$ 的全变差；见\cref{subsection:ch01-全变差}）。若\eqref{align:intro-0.4}成立，则对每个 $\mu\in\mathscr{R}_f(\mathbb{R}^N)$，存在 Marcinkiewicz 空间 $M^{N/(N-2)}(\mathbb{R}^N)$ 中唯一函数 $u$，在如下意义下解方程\eqref{align:intro-0.1}：即 $v:=\Delta u+\mu\in L^1(\mathbb{R}^N)$，且 $v=|u|^{p-1}u$ 在 $\mathbb{R}^N$ 中几乎处处成立（见文献~\cite{ref14} 附录 A）。反之，若条件\eqref{align:intro-0.4}不成立，则不存在任何函数使\eqref{align:intro-0.1}在 $\mu=\delta_{x_0}$（$x_0\in\mathbb{R}^N$）时成立（这可由一个直接的简单论证得到；见文献~\cite{ref14} 注记 A.4）。

上述结果如今已被充分理解。事实上，用于证明\eqref{align:intro-0.1}解存在性的方法是自然且被广泛使用的（可能有某些变体；例如见文献~\cite{ref27}）。考虑逼近问题族
\begin{align}
-\Delta u_n + |u_n|^{p-1}u_n = \mu_n \quad \text{在 }\mathbb{R}^N, \label{align:intro-0.5}
\end{align}
其中 $\{\mu_n\}\subseteq L^1(\mathbb{R}^N)\cap C^\infty(\mathbb{R}^N)$ 是 $\mu$ 的逼近序列，满足
\begin{align}
\|\mu_n\|_{L^1(\mathbb{R}^N)} \leqslant \|\mu\|_{\mathscr{R}_f(\mathbb{R}^N)}, \label{align:intro-0.6}
\end{align}
\begin{align}
\lim_{n\to\infty}\int_{\mathbb{R}^N}\rho\,\mu_n\,\mathrm{d}x
= \int_{\mathbb{R}^N}\rho\,\mathrm{d}\mu
\quad \text{对任意 }\rho\in C_0(\mathbb{R}^N) \label{align:intro-0.7}
\end{align}
（即 $\lim_{n\to\infty}\mu_n=\mu$ 在 $\mathscr{R}_f(\mathbb{R}^N)$ 中弱$^*$收敛；见\cref{section:ch05-Radon弱星}）。由文献~\cite{ref15}的结果，对每个 $n\in\mathbb{N}$，\eqref{align:intro-0.5}存在唯一解 $u_n\in M^{N/(N-2)}(\mathbb{R}^N)$。$u_n$ 的先验估计与一致界\eqref{align:intro-0.6}保证：(i) 序列 $\{u_n\}$ 在 $L^1_{\mathrm{loc}}(\mathbb{R}^N)$ 中相对紧；(ii) 序列 $\{v_n\}$，$v_n:=\Delta u_n+\mu_n$，在 $L^1(\mathbb{R}^N)$ 中有界。进而，若\eqref{align:intro-0.4}成立，则 $\{v_n\}$ 在 $\mathbb{R}^N$ 的有界子集上一致可积，因而由 Dunford--Pettis 定理在 $L^1_{\mathrm{loc}}(\mathbb{R}^N)$ 中相对弱紧（见文献~\cite{ref14} 引理 A.1 及\cref{theorem:ch02-dunfordpettis}）。于是在\eqref{align:intro-0.5}中令 $n\to\infty$ 并利用\eqref{align:intro-0.7}，即得所述存在性结果。

若\eqref{align:intro-0.4}不成立，则 $\{v_n\}$ 的一致可积性失效。由 Dunford--Pettis 定理，此时 $\{v_n\}$ 在 $L^1_{\mathrm{loc}}(\mathbb{R}^N)$ 中不相对弱紧，这与上述不存在性论断一致。然而即使在此情形，序列 $\{v_n\}$ 在 $L^1(\mathbb{R}^N)\subseteq\mathscr{R}_f(\mathbb{R}^N)$ 中仍有界，因而在 $\mathscr{R}_f(\mathbb{R}^N)$ 的弱$^*$拓扑下相对紧（见\cref{theorem:ch05-弱星紧性}）。故存在有限 Radon 测度 $v$，使得
\[
\lim_{n\to\infty}\int_{\mathbb{R}^N}\rho\,v_n\,\mathrm{d}x
= \int_{\mathbb{R}^N}\rho\,\mathrm{d}v
\quad \text{对任意 }\rho\in C_0(\mathbb{R}^N).
\]
这提示我们应在更大的 Radon 测度值解类中寻求方程\eqref{align:intro-0.1}（经适当修改后的版本）的解的存在性（见文献~\cite{ref97}）。

另一方面，使方程\eqref{align:intro-0.1}存在解 $u\in L^p(\mathbb{R}^N)$ 的那些测度 $\mu\in\mathscr{R}_f(\mathbb{R}^N)$ 已在文献~\cite{ref10}中刻画（另见文献~\cite{ref52}）：对任意 $p\in(1,\infty)$，这当且仅当 $\mu\in L^1(\mathbb{R}^N)+W^{-2,p}(\mathbb{R}^N)$ 时发生。由\cref{theorem:ch03-Lp可解性}，一个等价表述是：\eqref{align:intro-0.1}存在解 $u\in L^p(\mathbb{R}^N)$ 当且仅当 $\mu$ 关于 Sobolev 容量 $C_{2,q}$ 是弥散的，其中 $q:=p/(p-1)$ 为 $p$ 的共轭指数，亦即当且仅当
\begin{align}
\mu(E)=0 \quad \text{对任意满足 } C_{2,q}(E)=0 \text{ 的 Borel 集 }E\subseteq\mathbb{R}^N \label{align:intro-0.8}
\end{align}
（见\cref{subsection:ch01-Borel测度}与\cref{subsection:ch03-容量}）。

由\cref{remark:ch03-无零容量集}(ii)，若 $p\in(1,\infty)$ 且 $2q>N\Leftrightarrow p<N/(N-2)$，则不存在 Borel 集 $E\subseteq\mathbb{R}^N$ 使 $C_{2,q}(E)=0$，因而当\eqref{align:intro-0.4}成立时条件\eqref{align:intro-0.8}是空洞的（见\cref{proposition:ch03-点容量}）。反之，当 $p\geqslant N/(N-2)$ 时，对任意 $x_0\in\mathbb{R}^N$ 有 $C_{2,q}(\{x_0\})=0$，故 Dirac 质量 $\delta_{x_0}$ 关于 $C_{2,q}$ 是集中的，条件\eqref{align:intro-0.8}不满足（见\cref{definition:ch01-集中弥散}与\cref{proposition:ch03-单点容量}）。换言之，若 $p\geqslant N/(N-2)$，则 Dirac 质量对\eqref{align:intro-0.1}而言不是可去奇点（见文献~\cite{ref105}及其引文）。这就解释了上述关于方程\eqref{align:intro-0.1}的存在与不存在结果。

\begin{remark}\label{remark:intro-0.1.1}
对应于方程\eqref{align:intro-0.2}，由上述讨论可知：若 $N\geqslant 3$、$\mu$ 关于 Sobolev 容量 $C_{2,1/(1-\alpha)}$ 是集中的，且 $\alpha\leqslant (N-2)/N$，则不存在任何函数满足
\[
\lambda v - \Delta(|v|^{\alpha-1}v) = \mu \quad \text{在 }\mathbb{R}^N\ (\lambda,\alpha>0).
\]
这与文献~\cite{ref26}及文献~\cite{ref80} 定理 1、2 中的抛物结果一致（见下文）。
\end{remark}

需要提及，上述关于方程\eqref{align:intro-0.1}的存在与不存在结果刺激了多个方向的后续工作，特别是问题(E0)中不同的非线性椭圆算子，以及/或者(BC)中 $\partial\Omega\neq\emptyset$ 且 $g\in\mathscr{R}_f(\partial\Omega)$ 的情形（见文献~\cite{ref14,ref13,ref37,ref55,ref69,ref106}及其中引文）。

\subsubsection*{抛物情形}\label{subsubsection:intro-parabolic}

考虑快扩散多孔介质方程的 Cauchy 问题
\begin{align*}
\begin{cases}
\partial_t u = \Delta u^\alpha & \text{在 }\mathbb{R}^N\times(0,T)=:S,\\
u=u_0 & \text{在 }\mathbb{R}^N\times\{0\},
\end{cases}
\qquad(C)
\end{align*}
其中 $u_0$ 属于 $\mathbb{R}^N$ 上非负有限 Radon 测度的锥 $\mathscr{R}^+_f(\mathbb{R}^N)$（$N\geqslant 3$），$\alpha\in(0,(N-2)/N]$。文献~\cite{ref80}中，问题(C)的解按定义为任意 $u\in L^1(S)$，其在分布意义下满足第一式，并在窄收敛意义下满足第二式，即
\[
\operatorname*{ess\,lim}_{t\to 0^+}\int_{\mathbb{R}^N}\rho\,u(\cdot,t)\,\mathrm{d}x
= \int_{\mathbb{R}^N}\rho\,\mathrm{d}u_0
\quad \text{对任意 }\rho\in C_b(\mathbb{R}^N)
\]
（见\cref{definition:ch05-窄收敛}；此处 $C_b(\mathbb{R}^N)$ 表示 $\mathbb{R}^N$ 上连续有界实值函数空间）。文献证明：这样的解——更确切地说，解 $u\in C(0,T;L^1(\mathbb{R}^N))$——存在当且仅当 $u_0$ 关于 Sobolev 容量 $C_{2,1/(1-\alpha)}(\mathbb{R}^N)$ 是弥散的（见文献~\cite{ref80} 定理 1、2）。这与文献~\cite{ref26}一致，其中观察到：当 $u_0=\delta_0$ 时，问题(C)的这种解存在当且仅当 $\alpha\in((N-2)/N,1)$；事实上，对点 $x_0\in\mathbb{R}^N$，有 $C_{2,1/(1-\alpha)}(\{x_0\})=0$，因而 Dirac 质量 $\delta_{x_0}$ 关于 $C_{2,1/(1-\alpha)}$ 是集中的，当且仅当 $\alpha\in(0,(N-2)/N]$（见\cref{proposition:ch03-点容量}）。该结果也与\cref{theorem:ch6-正则化}（见等式\eqref{align:ch6-6.135}）一致。它与上述椭圆结果之间的联系是显然的（见上述注记）。

继问题(C)之后，初值为 Radon 测度的拟线性抛物方程初值问题已被广泛研究（特别见文献~\cite{ref16,ref26,ref35,ref36,ref81,ref104}及其中引文），且始终寻求函数值解。为凸显这一点，考虑初边值问题
\begin{align*}
\begin{cases}
\partial_t u = \Delta\phi(u) & \text{在 }\Omega\times(0,T)=:Q,\\
u=0 & \text{在 }\partial\Omega\times(0,T),\\
u=u_0 & \text{在 }\Omega\times\{0\},
\end{cases}
\qquad(P)
\end{align*}
这将在\cref{chapter:ch06}中详细讨论；此处 $\Omega\subseteq\mathbb{R}^N$ 为具有光滑边界 $\partial\Omega$ 的有界开集，$T>0$，$u_0\in\mathscr{R}_f(\Omega)$，并设 $\phi$ 在 $\mathbb{R}$ 上连续且非减。按我们的用语，问题(P)的解 $u$ 称为函数值的，若对任意 $t>0$ 有 $u(\cdot,t)\in L^1(\Omega)$。这种解可视为一类特殊的 Radon 测度值解，即问题(P)的这样的解 $u$：对正的时刻有 $u(\cdot,t)\in\mathscr{R}_f(\Omega)$（见文献~\cite{ref85,ref86,ref87}及其中引文）。这样的 Radon 测度值解是否在某个 $t>0$ 时变为函数值解，是一个正则性问题，也可视为奇点可去问题。不出所料，如下文所见，这一问题与所论发展方程特有的某种特征容量有关。

定义问题(P)的 Radon 测度值解，引出了如何定义这类测度的非线性函数的问题。文献~\cite{ref11}针对一个类似问题，通过启发式论证给出了一个自然定义，与文献~\cite{ref38,ref39}中给出的测度非线性函数概念一致。据此，我们把问题(P)的 Radon 测度值解理解为从 $(0,T)$ 到 $\mathscr{R}_f(\Omega)$ 的映射，并在如下弱意义下满足(P)：
\[
\int_0^T \langle u(\cdot,t),\partial_t\zeta(\cdot,t)\rangle\,\mathrm{d}t
= \iint_Q \nabla[\phi(u_r)]\cdot\nabla\zeta\,\mathrm{d}x\,\mathrm{d}t
- \langle u_0,\zeta(\cdot,0)\rangle
\]
（对一组合适的试验函数 $\zeta$；技术细节见\cref{chapter:ch06}）。特别地，要求测度 $u\in\mathscr{R}_f(Q)$ 属于空间 $L^\infty_{w^*}(0,T;\mathscr{R}_f(\Omega))$，因而对几乎所有 $t\in(0,T)$，$\Omega$ 上的 Radon 测度 $u(\cdot,t)$ 有定义；并要求函数 $\phi(u_r)$ 属于空间 $L^1(0,T;W^{1,1}_0(\Omega))$。此处 $u_r\in L^1(Q)$ 是 $u$ 的绝对连续部分 $u_{\mathrm{ac}}$ 关于 Lebesgue 测度的密度，$\langle\cdot,\cdot\rangle$ 表示 $C_0(\Omega)$ 与 $\mathscr{R}_f(\Omega)$ 之间的对偶映射（见\cref{chapter:ch02}与\cref{chapter:ch04}；按记号滥用，我们有时进一步把 $u_r$ 等同于 $u_{\mathrm{ac}}$）。

问题(P)的 Radon 测度值解可按前文椭圆情形所述方式构造：用合适序列 $\{u_{0n}\}\subseteq L^1(\Omega)$ 逼近初值 $u_0$，并研究相应逼近问题解序列在 $L^\infty_{w^*}(0,T;\mathscr{R}_f(\Omega))$ 中的收敛性（见\cref{chapter:ch06}）。就正则性而言，若当 $|z|\to\infty$ 时函数 $|\phi(z)|$ 的发散快于 $|z|^{(N-2)/N}$，则所构造的(P)的每个解都属于 $L^\infty(0,T;L^1(\Omega))$，即发生瞬时的 $\mathscr{R}_f$--$L^1$ 正则化（见\cref{proposition:ch6-瞬时正则}）。反之，若存在 $M>0$ 与 $\alpha\in(0,(N-2)/N]$ 使
\[
|\phi(z)| \leqslant M(1+|z|)^\alpha \quad (z\in\mathbb{R}),
\]
则问题(P)任意弱解的 $C_{2,1/(1-\alpha)}$ 集中部分在时间上为常数，这与文献~\cite{ref80}一致（见\cref{theorem:ch6-正则化}）。因此在此情形，每个这种解对任意正时刻都是 Radon 测度值的。

\subsection*{双曲守恒律的 Young 与 Radon 测度值解}\label{subsection:intro-hyperbolic}

考虑 Cauchy 问题
\begin{align*}
\begin{cases}
\partial_t u + \operatorname{div}[\phi(u)] = 0 & \text{在 }S,\\
u=u_0 & \text{在 }\mathbb{R}^N\times\{0\},
\end{cases}
\qquad(CL)
\end{align*}
其中 $\phi\equiv(\phi_1,\dots,\phi_N)\in(C^1(\mathbb{R}))^N$，$\phi_j:\mathbb{R}\to\mathbb{R}$ 非线性（$j=1,\dots,N$）。

\subsubsection*{Young 测度值解}\label{subsubsection:intro-young}

设 $u_0\in L^\infty(\mathbb{R}^N)$。证明(CL)弱解存在性的经典方法是粘性消失法，它依赖于相伴的抛物问题
\begin{align*}
\begin{cases}
\partial_t u_\epsilon + \operatorname{div}[\phi(u_\epsilon)] = \epsilon\,\Delta u_\epsilon & \text{在 }S,\\
u_\epsilon = u_{0\epsilon} & \text{在 }\mathbb{R}^N\times\{0\},
\end{cases}
\quad (\epsilon>0), \qquad(CL\epsilon)
\end{align*}
其中 $u_{0\epsilon}\in C_b(\mathbb{R}^N)$，且 $\|u_{0\epsilon}\|_{L^\infty(\mathbb{R}^N)}\leqslant\|u_0\|_{L^\infty(\mathbb{R}^N)}$。对每个 $\epsilon>0$，(CL$\epsilon$)存在唯一经典解 $u_\epsilon$，由极值原理满足不等式
\begin{align}
\|u_\epsilon\|_{L^\infty(S)} \leqslant \|u_0\|_{L^\infty(\mathbb{R}^N)}. \label{align:intro-0.9}
\end{align}
由不等式\eqref{align:intro-0.9}与 Banach--Alaoglu 定理，存在序列 $\{u_k\}\equiv\{u_{\epsilon_k}\}\subseteq L^\infty(S)$，$\epsilon_k\to0^+$（$k\to\infty$），以及 $u\in L^\infty(S)$，使 $\lim_{k\to\infty}u_k=u$ 在 $L^\infty(S)$ 中弱$^*$收敛。于是 $u$ 自然是(CL)弱解的候选。然而在(CL$\epsilon_k$)弱形式的右端令 $k\to\infty$，即
\begin{align}
\iint_S \{u_k\partial_t\zeta + \epsilon_k u_k\Delta\zeta\}\,\mathrm{d}x\,\mathrm{d}t
+ \int_{\mathbb{R}}u_0\zeta(x,0)\,\mathrm{d}x
= -\iint_S [\phi(u_k)]\cdot\nabla\zeta\,\mathrm{d}x\,\mathrm{d}t, \label{align:intro-0.10}
\end{align}
却颇为棘手，因为 $u_k$ 在 $L^\infty(S)$ 中弱$^*$收敛于 $u$，并不蕴含 $\phi_j(u_k)$ 弱$^*$收敛于 $\phi_j(u)$，$j=1,\dots,N$（见\cref{example:ch02-非线性弱星}(i)；此处 $\zeta$ 为合适的试验函数）。

为克服这一困难，我们需要更多信息，例如在 $S$ 上 $u_k\to u$ 几乎处处收敛（指 $\{u_k\}$ 的某个子列，为简便不重新标号）。事实上，由不等式\eqref{align:intro-0.9}、$\phi$ 的连续性与控制收敛定理，这一收敛将蕴含
\[
\lim_{k\to\infty}\iint_S[\phi(u_k)]\cdot\nabla\zeta\,\mathrm{d}x\,\mathrm{d}t
= \iint_S[\phi(u)]\cdot\nabla\zeta\,\mathrm{d}x\,\mathrm{d}t.
\]
于是在\eqref{align:intro-0.10}中令 $k\to\infty$，正如预期将得到
\begin{align}
\iint_S\{u\partial_t\zeta + [\phi(u)]\cdot\nabla\zeta\}\,\mathrm{d}x\,\mathrm{d}t
+ \int_{\mathbb{R}}u_0\zeta(x,0)\,\mathrm{d}x = 0. \label{align:intro-0.11}
\end{align}

为建立 $\{u_k\}$ 在 $S$ 上的几乎处处收敛，典型做法是置 $\mathbb{R}^N=\bigcup_{n=1}^\infty K_n$，其中每个 $K_n$ 紧，再用 Fréchet--Kolmogorov 定理证明：对任意固定的 $n\in\mathbb{N}$，限制序列 $\{u_k|_{K_n\times(0,T)}\}$ 在 $L^1(K_n\times(0,T))$ 中相对紧。然后由经典结果与对角论证即得结论（见文献~\cite{ref68} 第 2 章定理 4.62 的证明）。

下面概述另一种利用 Young 测度的方法（例如见文献~\cite{ref68}及其中引文）。记 $\lambda_N$ 为 $\mathbb{R}^N$ 上的 Lebesgue 测度，设 $U\subseteq S\subseteq\mathbb{R}^{N+1}$ 为有界开集，$\mathscr{R}^+_f(U\times\mathbb{R})$ 为 $U\times\mathbb{R}$ 上有限（正）Radon 测度的锥。按定义，$U\times\mathbb{R}$ 上的 Young 测度是任意满足
\[
\nu(E\times\mathbb{R}) = \lambda_{N+1}(E)
\quad \text{对任意 Borel 集 }E\subseteq U
\]
的 $\nu\in\mathscr{R}^+_f(U\times\mathbb{R})$。若 $f:U\to\mathbb{R}$ 可测，则满足
\[
\nu_f(E\times F) = \lambda_{N+1}(E\cap f^{-1}(F))
\]
对所有 Borel 集 $E\subseteq U$、$F\subseteq\mathbb{R}$ 成立的 Young 测度 $\nu_f$，称为与函数 $f$ 相伴的 Young 测度（见\cref{definition:ch02-Young测度}）。

记 $\mathcal{Y}(U;\mathbb{R})$ 为 $U\times\mathbb{R}$ 上 Young 测度的集合（$U\subseteq S$），$\mathcal{C}_b(U\times\mathbb{R})$ 为 $U\times\mathbb{R}$ 上有界 Carathéodory 被积函数的集合（见\cref{definition:ch05-Caratheodory}）。以下结果是\cref{proposition:ch05-离散分解}、\cref{theorem:ch05-弱紧Young}与\cref{proposition:ch05-biting}的推论。

\begin{proposition}\label{proposition:intro-0.2.1}
设 $\nu\in\mathcal{Y}(U;\mathbb{R})$。则对几乎所有 $(x,t)\in U$，存在 $\mathbb{R}$ 上的概率测度 $\nu_{(x,t)}$，使得对任意 $g\in C_b(U;\mathbb{R})$，映射 $(x,t)\mapsto\int_{\mathbb{R}}g(x,t,y)\,\mathrm{d}\nu_{(x,t)}(y)$ 在 $U$ 上可积，且
\[
\int_{U\times\mathbb{R}}g\,\mathrm{d}\nu
= \int_U \mathrm{d}x\,\mathrm{d}t\int_{\mathbb{R}}g(x,t,y)\,\mathrm{d}\nu_{(x,t)}(y).
\]
族 $\{\nu_{(x,t)}\}\equiv\{\nu_{(x,t)}\}_{(x,t)\in U}$ 是 $\mathbb{R}$ 上的参数化测度，称为 $\nu$ 的离散分解（见\cref{definition:ch05-参数测度}与\cref{definition:ch05-离散分解定义}）。由
\begin{align}
u^*(x,t) := \int_{\mathbb{R}}y\,\mathrm{d}\nu_{(x,t)}(y) \quad ((x,t)\in U) \label{align:intro-0.12}
\end{align}
定义的函数 $u^*$ 称为离散分解 $\{\nu_{(x,t)}\}$ 的重心。
\end{proposition}

\begin{theorem}\label{theorem:intro-0.2.2}
设 $\{u_n\}\subseteq L^1(U)$ 有界，$\{\nu_n\}\equiv\{\nu_{u_n}\}$ 为相伴 Young 测度序列。则：
\begin{enumerate}[(i)]
\item\label{item:intro-0.2.2-i} 存在子列 $\{u_{n_k}\}\subseteq\{u_n\}$ 与 $U\times\mathbb{R}$ 上的 Young 测度 $\nu$，使 $\nu_{n_k}\to\nu$ 窄收敛；
\item\label{item:intro-0.2.2-ii} 对任意 $f\in C(\mathbb{R})$ 使序列 $\{f(u_{n_k})\}\subseteq L^1(U)$ 有界且一致可积，有
\[
\int_U\mathrm{d}x\,\mathrm{d}t\int_{\mathbb{R}}|f(y)|\,\mathrm{d}\nu_{(x,t)}(y)<\infty,
\qquad
f(u_{n_k})\rightharpoonup f^* \quad \text{在 }L^1(U),
\]
其中
\begin{align}
f^*(x,t) := \int_{\mathbb{R}}f(y)\,\mathrm{d}\nu_{(x,t)}(y)
\quad \text{对几乎所有 }(x,t)\in U; \label{align:intro-0.13}
\end{align}
\item\label{item:intro-0.2.2-iii} 对任意 Carathéodory 函数 $g:U\times\mathbb{R}\to\mathbb{R}$ 使序列 $\{g(\cdot,u_{n_k})\}\subseteq L^1(U)$ 有界且一致可积，有
\[
\lim_{k\to\infty}\int_U g(x,t,u_{n_k}(x,t))\,\mathrm{d}x\,\mathrm{d}t
= \int_{U\times\mathbb{R}}g\,\mathrm{d}\nu;
\]
\item\label{item:intro-0.2.2-iv} [Biting 引理] 存在子列 $\{u_{n_{k_j}}\}\equiv\{u_{n_j}\}\subseteq\{u_{n_k}\}$ 与可测集序列 $\{U_j\}$，使对任意 $j\in\mathbb{N}$ 有 $U_{j+1}\subseteq U_j\subseteq U$，$\lambda_{N+1}(U_j)\to0$（$j\to\infty$），且序列 $\{u_{n_j}\chi_{U\setminus U_j}\}$ 一致可积。进而 $u^*\in L^1(U)$，且 $u_{n_j}\chi_{U\setminus U_j}\rightharpoonup u^*$ 在 $L^1(U)$。
\end{enumerate}
\end{theorem}

回到在\eqref{align:intro-0.10}中令 $k\to\infty$ 的问题。再次置 $\mathbb{R}^N=\bigcup_{n=1}^\infty K_n$，每个 $K_n$ 紧。对任意固定的 $n\in\mathbb{N}$，考虑限制序列 $\{u_k|_{K_n\times(0,T)}\}$，它在 $L^1(K_n\times(0,T))$ 中有界且由\cref{lemma:ch02-有界一致可积}一致可积。显然，序列 $\{\phi(u_k)|_{K_n\times(0,T)}\}$ 亦然。于是由\cref{theorem:intro-0.2.2}(ii)（另见\cref{remark:ch05-收敛注}(ii)）与对角论证，存在 $\{u_k\}$ 的子列（为简便不重新标号）与 Young 测度 $\nu\in\mathcal{Y}(S;\mathbb{R})$，使对任意有界开集 $U\subseteq S$，
\[
u_k\rightharpoonup u^*,\qquad
\phi_j(u_k)\rightharpoonup\phi_j^* \quad \text{在 }L^1(U),
\]
其中
\begin{align}
\phi_j^*(x,t) := \int_{\mathbb{R}}\phi_j(y)\,\mathrm{d}\nu_{(x,t)}(y)
\quad \text{对几乎所有 }(x,t)\in U\ (j=1,\dots,N). \label{align:intro-0.14}
\end{align}
由于 $\lim_{k\to\infty}u_k=u$ 在 $L^\infty(S)$ 中弱$^*$收敛，故 $u=u^*$。我们称偶 $(u,\nu)$ 为(CL)的弱 Young 测度值解，意指
\begin{align}
u(x,t) = \int_{\mathbb{R}}y\,\mathrm{d}\nu_{(x,t)}(y)
\quad \text{对几乎所有 }(x,t)\in S, \label{align:intro-0.15a}
\end{align}
且对任意 $\zeta\in C^1([0,T];C^1_c(\mathbb{R}))$，$\zeta(\cdot,T)=0$ 在 $\mathbb{R}$ 上，有
\begin{align}
\iint_S\{u\partial_t\zeta + \phi^*\cdot\nabla\zeta\}\,\mathrm{d}x\,\mathrm{d}t
+ \int_{\mathbb{R}}u_0\zeta(x,0)\,\mathrm{d}x = 0 \label{align:intro-0.15b}
\end{align}
其中 $\phi^*\equiv(\phi^*_1,\dots,\phi^*_N)$ 由\eqref{align:intro-0.14}给出。

至此我们仅用到了不等式\eqref{align:intro-0.9}。若流的每个分量 $\phi_j$ 都是真非线性的，则可说得更多（见\cref{chapter:ch07}中的假设 (A2)）。在此情形可证明 $\nu$ 实际上就是与 $u$ 相伴的 Young 测度，即 $\nu=\nu_u$。于是由\cref{proposition:ch05-Dirac测度}，对几乎所有 $(x,t)\in S$ 有 $\nu_{(x,t)}=\delta_{u(x,t)}$（见\eqref{align:ch05-5.53}），从而由\eqref{align:intro-0.14}有 $\phi^*=\phi(u)$ 在 $S$ 中几乎处处成立，等式\eqref{align:intro-0.15b}化归为\eqref{align:intro-0.11}。等价地，知道 $\nu=\nu_u$ 意味着（至多差一个子列）$u_k\to u$ 在 $S$ 上几乎处处（见\cref{proposition:ch05-几乎处处收敛}(ii)），因而我们可如前论证得到\eqref{align:intro-0.11}。

关键等式 $\nu=\nu_u$ 的证明依赖于所谓 Murat--Tartar 等式，而后者又用到 div--curl 引理（见文献~\cite{ref46} 推论 1.3.1 与定理 5.2.1，文献~\cite{ref68} 第 3.3 节；我们将用类似论证证明\cref{proposition:ch07-divcurl}）。

\subsubsection*{Radon 测度值解}\label{subsubsection:intro-radon1d}

一维空间中问题(CL)，
\begin{align*}
\begin{cases}
\partial_t u + \partial_x[\phi(u)] = 0 & \text{在 }\mathbb{R}\times(0,\infty),\\
u=u_0 & \text{在 }\mathbb{R}\times\{0\},
\end{cases}
\qquad(CL1)
\end{align*}
已在文献~\cite{ref67}中研究，其中 $u_0\in\mathscr{R}^+_f(\mathbb{R})$，且流 $\phi$ 在无穷远为超线性（模型情形为 $\phi(u)=u^m$，$m>1$）。(CL1)的熵解在如下意义下理解：
\begin{itemize}
\item 对正的时刻，$u$ 是函数：对每个 $\tau\in(0,T)$，$u\in L^\infty(0,T;L^1(\mathbb{R}))\cap L^\infty((\tau,T)\times\mathbb{R})$；
\item 对所有 $k\in\mathbb{R}$ 与 $\zeta\in C^1_c(S)$，$\zeta\geqslant0$，有熵不等式
\[
\iint_S\{|u-k|\partial_t\zeta + \operatorname{sgn}(u-k)[\phi(u)-\phi(k)]\partial_x\zeta\}\,\mathrm{d}x\,\mathrm{d}t \geqslant 0;
\]
\item 初值条件在窄收敛意义下满足：
\[
\operatorname*{ess\,lim}_{t\to0^+}\int_{\mathbb{R}}u(\cdot,t)\rho\,\mathrm{d}x
= \int_{\mathbb{R}}\rho\,\mathrm{d}u_0
\quad \text{对任意 }\rho\in C_b(\mathbb{R}).
\]
\end{itemize}

证明了如下结果。

\begin{proposition}\label{proposition:intro-0.2.3}
设 $\phi:[0,\infty)\to[0,\infty)$，$\phi(0)=0$ 为增函数，且对某个 $\alpha\in(0,1)$，$z\mapsto[\phi(z)]^{1-\alpha}$ 在 $(0,\infty)$ 上为凸。则(CL1)存在唯一熵解 $u\geqslant0$。
\end{proposition}

\cref{proposition:intro-0.2.3}的存在性部分由一个构造性过程给出，其中初值 Radon 测度 $u_0\geqslant0$ 用序列 $\{u_{0n}\}\subseteq L^\infty(\mathbb{R})\cap L^1(\mathbb{R})$ 逼近，$u_{0n}\to u_0$ 在 $\mathbb{R}$ 上窄收敛。若 $\phi(u)=u^m$，$m>1$，且 $u_0=\delta_0$，则该过程给出 Burgers 方程的基本解，称为 N 波：
\begin{align}
u(x,t) = \left(\frac{x}{mt}\right)^{1/(m-1)}\chi_A(x,t), \label{align:intro-0.16a}\\
A := \{(x,t)\mid 0\leqslant x\leqslant\left(\frac{m}{m-1}\right)^{(m-1)/m}t^{1/m},\ 0<t\leqslant T\}. \label{align:intro-0.16b}
\end{align}
因而此情形下(CL1)的解是函数值的（显然，$m\in(0,1)$ 时同样成立）。反之，在线性情形 $m=1$，初始奇点 $\delta_0$ 显然持续存在，故解对所有时刻都是 Radon 测度值的。

当研究具有界 $\phi$ 的(CL1)时会出现一种新情形；这一问题（由离子刻蚀的数学模型所推动，见文献~\cite{ref88}；另见文献~\cite{ref50}）在\cref{chapter:ch07}中处理。为说明其主要特征，考虑模型
\begin{align*}
\begin{cases}
\partial_t u + \partial_x\left[1-\dfrac{1}{(1+u)^m}\right] = 0 & \text{在 }\mathbb{R}\times(0,\infty),\ m>0,\\
u=\delta_0 & \text{在 }\mathbb{R}\times\{0\},
\end{cases}
\qquad(CL2)
\end{align*}
其中 $\phi(z)=1-\frac{1}{(1+z)^m}$ 为增且凹，并属于这样一类：对该类所构造的问题(CL2)的熵解是唯一的（见\cref{definition:ch07-流类}与\cref{theorem:ch07-存在唯一}）。如文献~\cite{ref67}那样，我们把(CL2)的熵解构造为逼近 Riemann 问题序列的熵解在（某合适拓扑下的）极限点：
\begin{align*}
\begin{cases}
\partial_t u_n + \partial_x\left[1-\dfrac{1}{(1+u_n)^m}\right] = 0 & \text{在 }S,\\
u_n = n\,\chi_{(-\frac{1}{2n},\frac{1}{2n})} & \text{在 }\mathbb{R}\times\{0\}.
\end{cases}
\qquad(Rn)
\end{align*}
对(Rn)的初等分析表明：一道激波从点 $(-\frac{1}{2n},0)$ 出发，一道稀疏波扇从 $(\frac{1}{2n},0)$ 出发，二者在点 $(x_n,t_n)$ 相遇，
\[
x_n := \frac{t_n}{2}\left[\frac{\phi(n)}{n}+\phi'(n)\right],
\qquad
t_n := \frac{1}{\phi(n)-n\phi'(n)} \quad (n\in\mathbb{N}),
\]
一道新的激波由此点出发。易见 $\lim_{n\to\infty}t_n=1$，$\lim_{n\to\infty}x_n=0$。这解释了以下结果（见文献~\cite{ref18}）。

\begin{proposition}\label{proposition:intro-0.2.4}
问题(CL2)的唯一（所构造的）熵解为 $u=u_r+u_s$，其中
\begin{align}
u_r(x,t) := \left[\left(\frac{mt}{x}\right)^{\frac{1}{1+m}}-1\right]\chi_A(x,t),\qquad
u_s(\cdot,t) := [1-t]_+\delta_0,
\nonumber\\
A := \{(x,t)\mid 0<x\leqslant mt,\ 0\leqslant t\leqslant1\}
\cup\{(x,t)\mid \xi(t)\leqslant x\leqslant mt,\ 1<t\leqslant T\}, \label{align:intro-0.17}
\end{align}
且 $\xi$ 求解问题
\[
\begin{cases}
\xi' = -\dfrac{(mt\xi-1)^{-\frac{m}{1+m}}-1}{(mt\xi-1)^{\frac{1}{1+m}}-1} & \text{在 }(1,T),\\
\xi(1,1)=0.
\end{cases}
\]
由\eqref{align:intro-0.17}，此情形下(CL1)解的奇异部分在等待时间 $t=1$ 之前非零，而当 $t\geqslant1$ 时消失，因而我们有延迟正则化。另一个重要注记是：对 $0<t<\tau_1=1$，奇异部分 $u_s(\cdot,t)$ 集中于 $x_0=0$（见\eqref{align:intro-0.17}）。当趋近于 $\operatorname{supp}u_s(\cdot,t)=\{0\}$ 时，对任意 $t\in(0,\tau_1)$，密度 $u_r(\cdot,t)$ 满足相容条件
\begin{align}
\lim_{x\to0^-}\phi(u_r(x,t))\equiv\phi(u_r(0^-,t))=0,
\qquad
\lim_{x\to0^+}\phi(u_r(x,t))\equiv\phi(u_r(0^+,t))=1. \label{align:intro-0.18}
\end{align}
我们请读者参阅\cref{chapter:ch07}，以领会在具有界 $\phi$ 的问题(CL1)唯一性证明中相容条件所起的本质作用（见\cref{definition:ch07-熵解}与\cref{theorem:ch07-唯一性}）。
\end{proposition}

\subsection*{不适定问题的 Young 与 Radon 测度值解}\label{subsection:intro-illposed}

现在考虑前向--后向抛物方程
\begin{align*}
\partial_t u = \Delta\phi(u) \quad \text{在 }\Omega\times(0,T)=:Q, \qquad(FB)
\end{align*}
其中 $\Omega\subseteq\mathbb{R}^N$（$N\geqslant1$）为开集，$\phi\in C^1(\mathbb{R})$ 非单调。方程(FB)可视为泛函
\begin{align}
\mathcal{E}[u] := \int_\Omega\Phi(u)\,\mathrm{d}x, \label{align:intro-0.19}
\end{align}
在 $H^{-1}$ 拓扑下的梯度流，其中 $\Phi'=\phi$。由于 $\phi$ 非单调，泛函 $\mathcal{E}$ 非凸，(FB)的初边值问题是不适定的。每当出现非凸泛函时——例如在相变、非线性弹性与图像处理中——与(FB)相关的动力学就有意义（见文献~\cite{ref29,ref73,ref79}）。不适定性表现为解的不唯一性（对满足下文假设 (A1)、(A2) 的 $\phi$，分别见文献~\cite{ref59}、文献~\cite{ref108}）。

在相变的 Landau 理论中，函数 $\Phi(u)=(1-u^2)^2$ 是双势阱位势，因而 $\phi$ 是三次的：
\begin{align*}
\begin{cases}
\phi'(u)>0 & \text{若 }u\in(-\infty,b)\cup(a,\infty),\\
\phi'(u)<0 & \text{若 }u\in(b,a);
\end{cases}
\qquad
\phi(u)\to\pm\infty \text{ 当 }u\to\pm\infty.
\qquad(A1)
\end{align*}
曲线 $v=\phi(u)$ 的三个单调分支记为
\begin{align}
\begin{aligned}
u &:= s_1(v),\ v\in(-\infty,B) \Leftrightarrow v=\phi(u),\ u\in(-\infty,b),\\
u &:= s_0(v),\ v\in(A,B) \Leftrightarrow v=\phi(u),\ u\in(b,a),\\
u &:= s_2(v),\ v\in(A,\infty) \Leftrightarrow v=\phi(u),\ u\in(a,\infty)
\end{aligned}
\label{align:intro-0.20}
\end{align}

\begin{figure}[ht]
\centering
\includegraphics[width=0.62\textwidth]{fig-cubic.png}
\caption{三次函数 $\phi$}\label{figure:intro-cubic}
\end{figure}

（此处 $A:=\phi(a)$，$B:=\phi(b)$；见\cref{figure:intro-cubic}）。启发式地，分支 $s_1$、$s_2$ 对应稳定相，$s_0$ 对应不稳定相，方程(FB)描述稳定相之间过渡的动力学。

不适定性促使人们对 Landau 理论作若干修正，特别包括非局部空间效应与/或时间延迟效应。在第一种情形，从所谓 Allen--Cahn 泛函
\[
\mathcal{E}[u] := \int_\Omega\left\{\Phi(u)+\frac{\kappa}{2}|\nabla u|^2\right\}\mathrm{d}x
\]
出发，得到 Cahn--Hilliard 方程
\begin{align}
\partial_t u = \Delta[\phi(u)-\kappa\Delta u] \quad (\kappa>0), \label{align:intro-0.21}
\end{align}
即(FB)的四阶正则化。反之，考虑时间延迟效应则导出(FB)的 Sobolev 正则化：
\[
\partial_t u = \Delta[\phi(u)+\epsilon\partial_t u] \quad (\epsilon>0).
\]
若同时纳入两种效应，则得到粘性 Cahn--Hilliard 方程
\[
\partial_t u = \Delta[\phi(u)-\kappa\Delta u+\epsilon\partial_t u].
\]
其他基于时间或空间离散的(FB)正则化已被多位作者使用~\cite{ref12,ref43,ref57,ref61}。

除相变外，研究方程(FB)的另一动机来自种群动力学，它在其中作为聚集种群离散模型的拟连续近似而出现（见文献~\cite{ref76,ref77}及其中引文）。在此框架中，$u\geqslant0$ 表示种群密度，$\phi\in C^2([0,\infty))$ 为具有如下性质的非单调函数：
\begin{align*}
\begin{cases}
\phi(u)>0 & \text{若 }u>0,\quad \phi(0)=0;\\
\phi'(u)>0 & \text{若 }0\leqslant u<\bar u,\quad \phi'(u)<0 \text{ 若 }u>\bar u;\\
\phi'(\bar u)=0,\ \phi''(\bar u)\neq0,\ \phi(u)\to0 \text{ 当 }u\to\infty;\\
\phi\in L^p(0,\infty) & \text{对某个 }p\in[1,\infty)
\end{cases}
\qquad(A2)
\end{align*}

\begin{figure}[ht]
\centering
\includegraphics[width=0.62\textwidth]{fig-perona.png}
\caption{Perona--Malik 型函数 $\phi$}\label{figure:intro-perona}
\end{figure}

（见\cref{figure:intro-perona}）。满足 (A2) 的函数常称为“Perona--Malik 型”，因为当 $N=1$ 时 Perona--Malik 方程写为
\begin{align}
\partial_t z = \partial_x[\phi(\partial_x z)] \quad \text{其中 }\phi(u)=\frac{u}{1+u^2}\ (u\geqslant0), \label{align:intro-0.22}
\end{align}
而上文的 $\phi$ 满足 (A2)（见文献~\cite{ref79}）。方程\eqref{align:intro-0.22}也刻画海洋中恒温（或盐度）层的形成~\cite{ref11}。在\eqref{align:intro-0.22}中置 $u:=\partial_x z$，得
\begin{align}
\partial_t u = \partial_{xx}\left(\frac{u}{1+u^2}\right) \quad (u\geqslant0), \label{align:intro-0.23}
\end{align}
这是(FB)的一个特例（需提及，$N=1$ 且 $\phi$ 为三次的方程(FB)已在文献~\cite{ref77}中研究）。注意\eqref{align:intro-0.22}中第一式是带 Perona--Malik 位势 $\Phi(u)=\log(1+u^2)$ 的\eqref{align:intro-0.19}的梯度流。

文献~\cite{ref11}给出了充分的建模论证，以说明\eqref{align:intro-0.22}的伪抛物正则化：
\[
\partial_t z = \partial_x[\phi(\partial_x z)] + \epsilon\partial_{tx}[\psi(\partial_x z)],
\]
其中带
\begin{align}
\psi(u) := -\phi(u) + \int_0^u\frac{\phi(z)}{z}\,\mathrm{d}z \quad (u\geqslant0) \label{align:intro-0.24}
\end{align}
的正则化项，是通过在(FB)中引入时间延迟效应而得到的。对\eqref{align:intro-0.22}中的函数 $\phi$，有 $\psi(u)=-\frac{u}{1+u^2}+\arctan u$，因而
\[
\psi'(u)>0 \quad \forall u\geqslant0, \qquad \psi(u)\to\frac{\pi}{2}\text{ 当 }u\to\infty.
\]
一般地，若 $\phi$ 满足 (A2)，则函数 $\psi:[0,\infty)\to[0,\infty)$ 满足下列要求：
\begin{align*}
\psi'(u)>0\ \forall u\geqslant0, \qquad \psi(u)\to\psi_\infty\in(0,\infty)\text{ 当 }u\to\infty. \qquad(A3)
\end{align*}

显然，伪抛物正则化弱于 Sobolev 正则化，后者形式上对应于 $\psi(u)=u$ 的情形。与\eqref{align:intro-0.22}和\eqref{align:intro-0.23}之间的对应一致，(FB)的伪抛物正则化写为
\[
\partial_t u = \Delta\phi(u) + \epsilon\partial_t\Delta[\psi(u)] \quad (u\geqslant0).
\]

\subsubsection*{Young 测度值解}\label{subsubsection:intro-youngfb}

考虑带 Sobolev 正则化的(FB)初边值问题：
\begin{align*}
\begin{cases}
\partial_t u = \Delta v & \text{在 }Q,\\
\partial_\nu v = 0 & \text{在 }\partial\Omega\times(0,T),\\
u=u_0 & \text{在 }\Omega\times\{0\},
\end{cases}
\qquad(S\epsilon)
\end{align*}
其中 $v:=\phi(u)+\epsilon\partial_t u$，$\Omega\subseteq\mathbb{R}^N$（$N\geqslant1$）为具有光滑边界 $\partial\Omega$ 的有界开集，$\partial_\nu$ 表示外法向导数。以下结果在文献~\cite{ref74}中证明（另见文献~\cite{ref76}）。

\begin{lemma}\label{lemma:intro-0.3.1}
设 $\phi\in C^1(\mathbb{R})$，$u_0\in C(\overline\Omega)$。则对每个 $\epsilon>0$，存在 $T_\epsilon>0$，使问题(S$\epsilon$)在 $Q_{T_\epsilon}:=\Omega\times(0,T_\epsilon)$ 中存在唯一强解 $(u_\epsilon,v_\epsilon)$。进而：
\begin{enumerate}[(i)]
\item\label{item:intro-0.3.1-i} 设 $g\in C^1(\mathbb{R})$，$g'\geqslant0$，$G(u):=\int_0^u g(\phi(s))\,\mathrm{d}s+c$（$c\in\mathbb{R}$）。则
\begin{align}
\int_\Omega G(u_\epsilon(x,t))\,\mathrm{d}x \leqslant \int_\Omega G(u_0(x))\,\mathrm{d}x
\quad \text{对任意 }t\in(0,T_\epsilon); \label{align:intro-0.25}
\end{align}
\item\label{item:intro-0.3.1-ii} 设 $u_1,u_2\in\mathbb{R}$，$u_1<u_2$，且
\begin{align}
\phi(u_1)\leqslant\phi(u)\leqslant\phi(u_2) \quad \text{对任意 }u\in[u_1,u_2]. \label{align:intro-0.26}
\end{align}
若对每个 $x\in\Omega$ 有 $u_0(x)\in[u_1,u_2]$，则对每个 $(x,t)\in Q_{T_\epsilon}$ 有 $u_\epsilon(x,t)\in[u_1,u_2]$。
\end{enumerate}
\end{lemma}

\begin{proof}
上述断言 (ii) 由 (i) 通过适当选取函数 $g$ 得到。为证 (i)，注意
\[
\begin{aligned}
\partial_t[G(u_\epsilon)]
&= g(\phi(u_\epsilon))\partial_t u_\epsilon
= g(v_\epsilon)\Delta v_\epsilon + [g(\phi(u_\epsilon))-g(v_\epsilon)]\Delta v_\epsilon\\
&= \operatorname{div}[g(v_\epsilon)\nabla v_\epsilon] - g'(v_\epsilon)|\nabla v_\epsilon|^2
+ [g(\phi(u_\epsilon))-g(v_\epsilon)]\frac{v_\epsilon-\phi(u_\epsilon)}{\epsilon}
\leqslant 0.
\end{aligned}
\]
将上式在 $\Omega$ 上积分即得\eqref{align:intro-0.25}。
\end{proof}

若 (A1) 成立，则对适当选取的 $u_1<0<u_2$（$|u_1|$、$|u_2|$ 充分大），函数 $\phi$ 满足\eqref{align:intro-0.26}。于是我们得到 $\{u_\epsilon\}$ 的一致 $L^\infty$ 估计。下文\eqref{align:intro-0.27}中的其他估计由不等式\eqref{align:intro-0.25}的证明类似得到，故有如下结果。

\begin{proposition}\label{proposition:intro-0.3.2}
设 (A1) 成立，$u_0\in C(\overline\Omega)$。则：
\begin{enumerate}[(i)]
\item\label{item:intro-0.3.2-i} 对每个 $\epsilon>0$，正则化问题(S$\epsilon$)存在唯一强解 $(u_\epsilon,v_\epsilon)$，其中 $v_\epsilon=\phi(u_\epsilon)+\epsilon\partial_t u_\epsilon$；
\item\label{item:intro-0.3.2-ii} 存在 $M>0$，使对任意 $\epsilon>0$，
\begin{align}
\max\left\{\|u_\epsilon\|_{L^\infty(Q)},\ \sqrt{\epsilon}\,\|\partial_t u_\epsilon\|_{L^2(Q)},\ \|v_\epsilon\|_{L^\infty(Q)},\ \|v_\epsilon\|_{L^2(0,T;H^1(\Omega))}\right\} \leqslant M. \label{align:intro-0.27}
\end{align}
\end{enumerate}
\end{proposition}

由不等式\eqref{align:intro-0.27}，存在序列 $\{\epsilon_k\}$（$\epsilon_k\to0^+$，$k\to\infty$）、$\{u_k\}\equiv\{u_{\epsilon_k}\}$、$\{v_k\}\equiv\{v_{\epsilon_k}\}$，以及 $u\in L^\infty(Q)$、$v\in L^\infty(Q)\cap L^2(0,T;H^1(\Omega))$，使
\begin{align}
\begin{cases}
u_k \stackrel{*}{\rightharpoonup} u & \text{在 }L^\infty(Q),\\
\phi(u_k) \stackrel{*}{\rightharpoonup} v & \text{在 }L^\infty(Q),\\
v_k \stackrel{*}{\rightharpoonup} v & \text{在 }L^\infty(Q),\\
v_k \rightharpoonup v & \text{在 }L^2(0,T;H^1(\Omega))
\end{cases}
\label{align:intro-0.28}
\end{align}
（注意序列 $\{v_k\}$ 与 $\{\phi(u_k)\}$ 在 $L^\infty(Q)$ 中有相同的弱$^*$极限）。由\eqref{align:intro-0.28}，在(S$\epsilon$)$_{\epsilon_k}$ 的弱形式中令 $k\to\infty$，易得
\begin{align}
\iint_Q\{u\partial_t\zeta - \nabla v\cdot\nabla\zeta\}\,\mathrm{d}x\,\mathrm{d}t
+ \int_\Omega u_0\zeta(x,0)\,\mathrm{d}x = 0 \label{align:intro-0.29}
\end{align}
对任意 $\zeta\in C^1([0,T];C^1(\overline\Omega))$、$\zeta(\cdot,T)=0$ 在 $\Omega$ 上成立。

由于\eqref{align:intro-0.28}中的收敛并不蕴含等式 $v=\phi(u)$，$u$ 不必是问题
\begin{align*}
\begin{cases}
\partial_t u = \Delta[\phi(u)] & \text{在 }Q,\\
\partial_\nu[\phi(u)] = 0 & \text{在 }\partial\Omega\times(0,T),\\
u=u_0 & \text{在 }\Omega\times\{0\}
\end{cases}
\qquad(S)
\end{align*}
的弱解。然而，(S)的弱 Young 测度值解却很容易得到。由于 $\|u_k\|_{L^\infty(Q)}\leqslant M$（见\eqref{align:intro-0.27}）且 $Q$ 有界，序列 $\{u_k\}$ 在 $L^1(Q)$ 中有界且一致可积（见\cref{lemma:ch02-有界一致可积}），序列 $\{\phi(u_k)\}$ 亦然。于是存在 $\{u_k\}$ 的子列（不重新标号）与 Young 测度 $\nu\in\mathcal{Y}(Q;\mathbb{R})$，使
\[
u_k\rightharpoonup u^*,\qquad \phi(u_k)\rightharpoonup\phi^* \quad \text{在 }L^1(Q),
\]
其中 $u^*$、$\phi^*$ 分别由\eqref{align:intro-0.12}、\eqref{align:intro-0.13}在 $U=Q$、$f=\phi$ 时给出。由\eqref{align:intro-0.28}与上述收敛得 $u=u^*$、$v=\phi^*$，从而由\eqref{align:intro-0.29}有
\begin{align}
\iint_Q\{u\partial_t\zeta - \nabla\phi^*\cdot\nabla\zeta\}\,\mathrm{d}x\,\mathrm{d}t
+ \int_{\mathbb{R}}u_0\zeta(x,0)\,\mathrm{d}x = 0. \label{align:intro-0.30}
\end{align}
因此偶 $(u,\nu)$ 是(S)的弱 Young 测度值解。

若 $\phi$ 满足假设 (A1)，则有如下 Young 测度 $\nu$ 的刻画~\cite{ref83}（另见文献~\cite{ref82,ref84}）：
\begin{align}
\nu_{(x,t)} = \sum_{i=0}^2 \lambda_i(x,t)\,\delta_{s_i(v(x,t))}
\quad \text{对几乎所有 }(x,t)\in Q. \label{align:intro-0.31}
\end{align}
此处系数 $\lambda_i$ 取值于 $[0,1]$，$\sum_{i=0}^2\lambda_i=1$，函数 $s_i(v)$ 由\eqref{align:intro-0.20}定义。由等式\eqref{align:intro-0.12}与\eqref{align:intro-0.31}有
\begin{align}
u(x,t) = \sum_{i=0}^2\lambda_i(x,t)\,s_i(v(x,t))
\quad \text{对几乎所有 }(x,t)\in Q. \label{align:intro-0.32}
\end{align}

\begin{remark}\label{remark:intro-0.3.1}
由\eqref{align:intro-0.32}，$u$ 是不同相的叠加。按定义，(S)的两相解是指上述这样的解：在 $Q$ 中几乎处处 $\lambda_0=0$，且在 $V_i$ 上几乎处处 $\lambda_i=1$（$i=1,2$），其中 $Q=V_1\cup V_2$，$V_1\cap V_2=\emptyset$，公共边界 $V_1\cap V_2$ 光滑。

当 $N=1$ 时，两相解的一个主要定性特征是表现出迟滞效应：唯一允许的相变是那些从三次曲线的极值点 $(b,B)$ 与 $(c,A)$ 出发发生的相变（见\cref{figure:intro-cubic}）。其证明与双曲守恒律分段光滑解的 Oleinik 熵条件证明类似（后者在某种程度上正是(S)两相解的对应物；见文献~\cite{ref48,ref54,ref70,ref71,ref82,ref83,ref84,ref98}）。

除其他点外，这一特征指出：由 Sobolev 正则化得到的(S)解，与由 Cahn--Hilliard 正则化得到的解确实不同。事实上，文献~\cite{ref12}证明，在一维环面上\eqref{align:intro-0.21}的解当 $\kappa\to0^+$ 时的极限动力学由 Maxwell 等面积律支配：
\[
\int_{-1}^1\phi^{**}(u)\,\mathrm{d}u = 0 = \int_{-1}^1\phi(u)\,\mathrm{d}u,
\]
这与迟滞现象不相容（此处 $\phi^{**}:=(\Phi^{**})'$，$\Phi^{**}$ 为双势阱位势 $\Phi$ 的凸包络）。文献~\cite{ref43}用隐式变分格式就一个梯度方程的不适定问题得到了类似结果。同类问题在文献~\cite{ref30,ref95}中也分别用 Sobolev 与 Cahn--Hilliard 正则化作了研究（见\cref{chapter:ch08}）。
\end{remark}

\subsubsection*{Radon 测度值解}\label{subsubsection:intro-radonfb}

设 $\phi$、$\psi$ 分别满足 (A2)、(A3)。由于当 $u\to\infty$ 时 $\psi'(u)\to0$，初边值问题
\begin{align*}
\begin{cases}
\partial_t u = \partial_{xx}[\phi(u)] + \epsilon\partial_{txx}[\psi(u)] & \text{在 }Q,\\
\phi(u) + \epsilon\partial_t[\psi(u)] = 0 & \text{在 }\partial\Omega\times(0,T),\\
u=u_0 & \text{在 }\Omega\times\{0\},\\
u\geqslant0,
\end{cases}
\qquad(PM\epsilon)
\end{align*}
可以想见，对 $\epsilon>0$ 所表现出的定性特征，与带 Sobolev 正则化的问题
\begin{align*}
\begin{cases}
\partial_t u = \partial_{xx}[\phi(u)] + \epsilon\partial_{txx}u & \text{在 }Q,\\
\phi(u) + \epsilon\partial_t u = 0 & \text{在 }\partial\Omega\times(0,T),\\
u=u_0 & \text{在 }\Omega\times\{0\},\\
u\geqslant0,
\end{cases}
\qquad(\Sigma\epsilon)
\end{align*}
当 $\epsilon\to0^+$ 时所表现的特征，就奇点性态而言是相同的。问题(PM$\epsilon$)（$u_0\in\mathscr{R}^+_f(\Omega)$）已在文献~\cite{ref99,ref100}中研究。其最显著特征（显然与伪抛物正则化的微弱性有关）是：当 $u_0$ 光滑时，正的时刻会自发出现奇点（见文献~\cite{ref22} 定理 4.2、\cref{definition:ch01-集中测度}与\cref{theorem:ch03-奇性}）：

\begin{proposition}\label{proposition:intro-0.3.3}
设 (A2)--(A3) 成立。存在(PM$\epsilon$)的解（$u_0\in C(\overline\Omega)$，$u_0\geqslant0$），使得对某个 $t\in(0,T)$，要么 $u_s(\cdot,t)$ 是纯原子的且含可数多个 Dirac 质量，要么 $u_s(\cdot,t)$ 是奇异连续的。
\end{proposition}

该结果也依赖于 $u$ 取大值时 $\phi'(u)$ 的负性，它导致集中现象。这与文献~\cite{ref11}的结果形式上一致：问题
\[
\begin{cases}
\partial_t z = \partial_x[\phi(\partial_x z)] + \epsilon\partial_{tx}[\psi(\partial_x z)] & \text{在 }Q,\\
\phi(\partial_x z) + \epsilon\partial_t[\psi(\partial_x z)] = 0 & \text{在 }\partial\Omega\times(0,T),\\
z=z_0\in C^1(\overline\Omega) & \text{在 }\Omega\times\{0\}
\end{cases}
\]
的解在正的时刻可能关于 $x$ 间断。

最后，比较在问题($\Sigma\epsilon$)中令 $\epsilon\to0^+$ 时 (A1) 与 (A2) 的效果是有趣的。若 $\phi$ 为三次，则每个充分大的区间 $[u_1,u_2]$ 都满足条件\eqref{align:intro-0.26}，故\cref{lemma:intro-0.3.1}(ii) 给出 $\{u_\epsilon\}$ 的一致 $L^\infty$ 估计。因此当 $\epsilon\to0^+$ 时，我们得到 $u\in L^\infty(Q)$ 的等式\eqref{align:intro-0.30}，故在此情形不出现奇点。反之，若 $\phi$ 为 Perona--Malik 型，则只有那些 $\phi'$ 为正的区间 $[u_1,u_2]\subseteq[0,\infty)$ 满足\eqref{align:intro-0.26}。因而 $\{u_\epsilon\}$ 的一致 $L^\infty$ 估计仅当 $\{u_\epsilon\}$ 取稳定相中的值时成立，而这是平凡的。

然而在后一情形，可证明半直线 $[0,\infty)$ 对(PM$\epsilon$)的解是正不变的（见文献~\cite{ref96}）。于是，若 $u_0\in L^1(\Omega)$，$u_0\geqslant0$，则质量守恒
\[
\int_\Omega u_\epsilon(x,t)\,\mathrm{d}x = \int_\Omega u_0(x)\,\mathrm{d}x \quad (t\in(0,T))
\]
给出 $\{u_\epsilon\}$ 在 $L^\infty(0,T;L^1(\Omega))$ 中的一致估计：
\[
\|u_\epsilon(\cdot,t)\|_{L^1(\Omega)}
= \int_\Omega u_\epsilon(x,t)\,\mathrm{d}x
= \int_\Omega u_0(x)\,\mathrm{d}x
= \|u_0\|_{L^1(\Omega)}.
\]
由此
\[
\|u_\epsilon\|_{L^1(Q)} \leqslant T\|u_0\|_{L^1(\Omega)} \quad \forall\epsilon>0.
\]
于是由 biting 引理（见\cref{theorem:intro-0.2.2}(iv)与\cref{theorem:ch05-biting定理}），存在序列 $\{u_k\}\equiv\{u_{\epsilon_k}\}\subseteq\{u_\epsilon\}$、Young 测度 $\nu\in\mathcal{Y}(Q;\mathbb{R})$、Radon 测度 $\sigma\in\mathscr{R}^+_f(Q)$，以及序列 $Q_{k+1}\subseteq Q_k\subseteq Q$（其 Lebesgue 测度 $\lambda_{N+1}(Q_k)$ 当 $\epsilon_k\to0^+$ 时趋于零），使得（可能至多差一个不重新标号的子列）
\[
u_k\chi_{Q\setminus Q_k}\rightharpoonup u^* := \int_{[0,\infty)}y\,\mathrm{d}\nu_{(x,t)}(y)
\quad \text{在 }L^1(Q),
\]
以及
\begin{align}
u_k\chi_{Q_k} \stackrel{*}{\rightharpoonup}\sigma \quad \text{在 }\mathscr{R}_f(Q). \label{align:intro-0.33}
\end{align}
如\cref{proposition:intro-0.3.2}那样，也易验证存在 $v\in L^\infty(Q)\cap L^2(0,T;H^1_0(\Omega))$，使
\[
v_k\stackrel{*}{\rightharpoonup}v \text{ 在 }L^\infty(Q),
\qquad
v_k\rightharpoonup v \text{ 在 }L^2(0,T;H^1_0(\Omega)),
\]
其中 $v_k:=\phi(u_k)+\epsilon_k\partial_t u_k$。于是在($\Sigma\epsilon$)$_{\epsilon_k}$ 的弱形式中令 $k\to\infty$，易得
\[
\iint_Q(u^*\partial_t\zeta - \nabla v\cdot\nabla\zeta)\,\mathrm{d}x\,\mathrm{d}t
+ \int_\Omega u_0\zeta(x,0)\,\mathrm{d}x
= -\langle\sigma,\partial_t\zeta\rangle
\]
对任意 $\zeta\in C^1([0,T];C^1_c(\Omega))$、$\zeta(\cdot,T)=0$ 在 $\Omega$ 上成立。上述等式右端（由\eqref{align:intro-0.33}中的收敛产生）刻画了($\Sigma\epsilon$)当 $\epsilon\to0^+$ 时可能的集中现象（在(PM$\epsilon$)中，因正则化更弱，$\epsilon>0$ 时已观察到这一点）。另一方面，如三次 $\phi$ 的情形，可验证 $v=\phi^*$，其中
\[
\phi^*(x,t):=\int_{\mathbb{R}}\phi(y)\,\mathrm{d}\nu_{(x,t)}(y)
\quad \text{对几乎所有 }(x,t)\in Q.
\]
因此，在 Perona--Malik 情形，需要同时用 Radon 测度与 Young 测度来描述($\Sigma\epsilon$)的粘性消失极限。

上述讨论提示出如下图景：
\begin{enumerate}[(i)]
\item $\phi$ 在无穷远的性态决定解的奇异部分关于 Lebesgue 测度的性态，而单调性的改变则影响绝对连续部分的性态；
\item 奇异部分的性态由 Radon 测度描述，并与奇点的消失、持续或出现（集中或正则化现象）有关；
\item 绝对连续部分的性态由 Young 测度描述，并与振荡现象有关。
\end{enumerate}

\section*{数学工具}\label{section:intro-tools}

让我们从上述讨论中得出一些结论。处理涉及测度的非线性微分问题时，总要用到某种逼近过程。逼近问题的解是某个函数空间中的一个族，而该函数空间自然嵌入有限 Radon 测度的 Banach 空间。当某个正则化参数趋于零时，该族在某合适拓扑下的极限点，便是原问题解的自然候选。

因此，证明解存在性的主要工具，由关于有限 Radon 测度序列在各种拓扑下的收敛性与紧性的结果所提供。这些结果适合（且其证明用到）定义在某给定 Banach 空间上的测度值函数的 Lebesgue 空间这一一般框架。确实，证明的清晰与优雅极大地受益于对泛函分析背景的清楚理解。

把一般理论应用于逼近问题解序列，便得到收敛性结果，使我们能在逼近问题的弱形式中取极限。视问题的性质而定，这类结果同时刻画集中现象与（若存在的话）振荡现象（见\cref{section:ch05-收敛}）。

在这一一般框架中，一个主要问题是 Radon 测度值解的正则性，因而是奇点可能的存在性、位置与演化。如前所述，这些特征与问题特有的合适容量有关。因此，容量理论也是研究这类问题的重要工具。

\section*{内容简述与阅读指引}\label{section:intro-guide}

本书第一部分致力于一般理论。第 1、2 章具有预备性质，分别介绍抽象测度论以及标量函数可测性与积分的一般结果。第 3、4、5 章是本书核心，分别处理容量理论、向量积分与向量函数的 Lebesgue 空间，以及有限 Radon 测度序列的收敛。

在第二部分（第 6、7、8 章）中，我们描述一般理论的三个应用。每个应用对应于本引言所提及的三个研究领域之一（拟线性抛物问题、双曲守恒律、不适定发展问题）。

各部分开头给出该部分章节内容的说明。必要的拓扑预备知识在书末附录 A 中回顾。参考文献虽远非穷尽，但对本书所论各主题给出了充分的引证。

补充几点实用说明：
\begin{enumerate}[1.]
\item 所谓“标量函数”，指任意实值函数；与“向量函数”相对，后者指取值于 $\mathbb{R}$ 以外任意 Banach 空间的映射。
\item 所考虑的向量空间均为实的。
\item 在每节内，定理、命题、引理与推论依次编号（例如定理 1.2.1 之后是命题 1.2.2 等），而定义、注记与例题单独编号。公式在每章内依次编号。
\end{enumerate}

\end{document}
